% \iffalse meta-comment
% Copyright (C) 2026 Andreas Willibald Weber
%
% This work may be distributed and/or modified under the conditions of the
% LaTeX Project Public License, either version 1.3c of this license or (at your
% option) any later version. The latest version of the license is available at
% https://www.latex-project.org/lppl.txt.
%
% This work has the LPPL maintenance status `maintained'. The Current
% Maintainer of this work is Andreas Willibald Weber.
%
% This file is part of the `kanban' package; the complete list of files
% constituting the work appears in the main source file kanban.dtx.
% \fi
%
%    \begin{macrocode}
%<*package>
%    \end{macrocode}
%
% \subsection{Card components and renderers}
%
% Metadata elements are small inline boxes. Their semantic colors are stable
% across modes; only the card surface, card text, and card border respond to
% the light or dark board mode. Public, namespaced styles are appended after
% structural defaults. Fixed and breakable layouts reuse the same card renderer
% but provide different column containers. Rendering is deliberately delayed:
% task declarations store calls to these components, and the environment
% executes those calls only after it knows all column counts and dimensions.
%
% \paragraph{Structural and extension styles.}
% The first style registry distinguishes fixed package structure from supported
% user extension points. Empty public styles such as |kanban/card| may be
% extended with |\tcbset|, while private styles establish the common geometry
% of inline pills and the stable author and deadline colors.
%
% |\tcbset| defines |tcolorbox| keys globally. A |/.style| value expands to a
% reusable set of keys. |on line| makes a box participate in the current text
% line; |boxsep| and the side-specific dimensions control padding; |arc| rounds
% corners; |boxrule=0pt| removes the frame; |nobeforeafter| suppresses paragraph
% material; and |box align=base| aligns the box with surrounding text. The
% private |kanban@chip| and |kanban@due| styles inherit |kanban@pill| and add
% semantic foreground and background macros. Because the public extension key
% is applied last in each component definition below, a document may override
% presentation without replacing structural collection logic.
%    \begin{macrocode}
\tcbset{
  kanban/header/.style={},
  kanban/card/.style={},
  kanban/column-body/.style={},
  kanban/priority-chip/.style={},
  kanban/author-chip/.style={},
  kanban/deadline-chip/.style={},
  kanban/completed-chip/.style={},
  kanban@pill/.style={on line,boxsep=0pt,left=1.2mm,right=1.2mm,
    top=.4mm,bottom=.4mm,arc=.7mm,boxrule=0pt,nobeforeafter,
    box align=base},
  kanban@chip/.style={kanban@pill,colback=\kanban@authorback,
    coltext=\kanban@authortext},
  kanban@due/.style={kanban@pill,colback=\kanban@deadlineback,
    coltext=\kanban@deadlinetext}}
%    \end{macrocode}
%
% \paragraph{Priority and author chips.}
% A priority chip maps a normalized semantic value to the board-local label,
% semantic background, and accessible foreground. An empty or unknown value
% produces no chip. The author component has no semantic classification and
% therefore renders any non-blank author text with the stable author style.
%
% Both commands are checked LaTeX2e definitions made by |\newcommand|.
% |\ifblank| avoids creating empty inline boxes. The priority renderer initializes
% three temporary macros and fills them through |\ifstrequal| comparisons;
% |\ifdefempty| then guards the actual box. |\tcbox| draws a one-line
% |tcolorbox|, and its optional key list combines the structural pill, the
% selected colors, and the public extension style. |\bfseries| and
% |\scriptsize| are declarations scoped by the box content. The author renderer
% uses the inherited |kanban@chip| colors and applies |kanban/author-chip| last.
%    \begin{macrocode}
\newcommand\kanban@prioritypill[1]{\ifblank{#1}{}{%
  \def\kb@pl{}\def\kb@pc{}\def\kb@pt{}%
  \ifstrequal{#1}{high}{\def\kb@pl{\kanban@boardhighlabel}%
    \def\kb@pc{kanbanHigh}\def\kb@pt{\kanban@hightext}}{}%
  \ifstrequal{#1}{medium}{\def\kb@pl{\kanban@boardmediumlabel}%
    \def\kb@pc{kanbanMed}\def\kb@pt{\kanban@mediumtext}}{}%
  \ifstrequal{#1}{med}{\def\kb@pl{\kanban@boardmediumlabel}%
    \def\kb@pc{kanbanMed}\def\kb@pt{\kanban@mediumtext}}{}%
  \ifstrequal{#1}{low}{\def\kb@pl{\kanban@boardlowlabel}%
    \def\kb@pc{kanbanLow}\def\kb@pt{\kanban@lowtext}}{}%
  \ifdefempty\kb@pl{}{\tcbox[kanban@pill,
    colback=\kb@pc,colframe=\kb@pc,coltext=\kb@pt,
    kanban/priority-chip]{\bfseries\scriptsize\kb@pl}}}}
\newcommand\kanban@authorchip[1]{\ifblank{#1}{}{%
  \tcbox[kanban@chip,kanban/author-chip]{\scriptsize #1}}}
%    \end{macrocode}
%
% \paragraph{The deadline clock symbol.}
% Deadline chips include a compact clock drawn from primitives so the package
% does not depend on an icon font. Its dimensions are fixed in points and are
% small enough to align with script-sized metadata text.
%
% The inline |\tikz| command creates a self-contained picture. Its |baseline|
% option lowers the picture relative to the surrounding baseline. The first
% |\draw| command creates a circle, and the next two draw the clock hands; each
% uses the same thin |line width|. Coordinates and lengths are local to this
% picture, so the symbol does not alter the dimensions or style of later TikZ
% material.
%    \begin{macrocode}
\newcommand\kb@clock{\tikz[baseline=-.55ex]{%
  \draw[line width=.3pt] (0,0) circle (1.7pt);
  \draw[line width=.3pt] (0,0)--(0,1.1pt);
  \draw[line width=.3pt] (0,0)--(1pt,.2pt);}}
%    \end{macrocode}
%
% \paragraph{Deadline and completion chips.}
% Deadline rendering first classifies the raw value and then prepares its
% display notation. A manually |late| task or an automatically overdue date
% receives the high semantic color; a date equal to the reference date receives
% the medium color; future dates, project labels, and descriptive text retain
% the neutral deadline style. Completion is a separate Boolean flag rendered
% with the low semantic color. These colors remain identical in light and dark
% board modes.
%
% |\ifboolexpr| combines the string test for the stored late flag with the
% |etoolbox| overdue Boolean. Its |test| wrapper embeds |\ifstrequal|, while
% |or bool| reads the named Boolean. The nested |\ifbool| selects due-now versus
% neutral output. Each branch invokes |\tcbox| with the public deadline style
% last and includes the clock, spacing command |\,|, size, and optional bold
% declaration. |\kanban@completedchip| uses |\ifstrequal| to recognize the
% stored flag and otherwise expands to nothing.
%    \begin{macrocode}
\newcommand\kanban@deadlinechip[2]{\ifblank{#2}{}{%
  \kb@evaloverdue{#2}%
  \kanban@preparedeadline{#2}%
  \ifboolexpr{test {\ifstrequal{#1}{1}} or bool {kb@overdue}}
    {\tcbox[kanban@due,colback=kanbanHigh,coltext=\kanban@hightext,
        kanban/deadline-chip]
      {\kb@clock\,\scriptsize\bfseries \kb@displaydeadline}}
    {\ifbool{kb@duenow}
      {\tcbox[kanban@due,colback=kanbanMed,coltext=\kanban@mediumtext,
          kanban/deadline-chip]
        {\kb@clock\,\scriptsize\bfseries \kb@displaydeadline}}
      {\tcbox[kanban@due,kanban/deadline-chip]
        {\kb@clock\,\scriptsize \kb@displaydeadline}}}}}
\newcommand\kanban@completedchip[1]{\ifstrequal{#1}{1}
  {\tcbox[kanban@pill,colback=kanbanLow,colframe=kanbanLow,
      coltext=\kanban@lowtext,kanban/completed-chip]
    {\bfseries\scriptsize\kanban@boarddonelabel}}{}}
%    \end{macrocode}
%
% \paragraph{Priority-value validation.}
% Task options must store only priority values understood by the renderer. The
% setter normalizes case and surrounding spaces, accepts the four documented
% spellings, treats a blank choice as omission, and warns when non-blank input
% is unknown. The normalized value, rather than the original spelling, is
% stored with the task.
%
% The implementation reuses |\kanban@normalizechoice|. It initializes the
% task-priority macro |\kbtask@prio| with |\def|. Nested |\ifdefstring| tests select the accepted
% replacement. The final |\ifdefempty| distinguishes an intentional blank from
% a rejected choice before issuing |\PackageWarning|. Because all definitions
% occur while a task is being collected inside an environment group, the
% temporary value cannot leak to later tasks.
%    \begin{macrocode}
\newcommand\kanban@setpriority[1]{%
  \kanban@normalizechoice{#1}%
  \def\kbtask@prio{}%
  \ifdefstring{\kb@choice}{high}{\def\kbtask@prio{high}}{%
    \ifdefstring{\kb@choice}{medium}{\def\kbtask@prio{medium}}{%
      \ifdefstring{\kb@choice}{med}{\def\kbtask@prio{med}}{%
        \ifdefstring{\kb@choice}{low}{\def\kbtask@prio{low}}{%
          \ifdefempty\kb@choice{}{\PackageWarning{kanban}
            {unknown priority '#1'; omitting the priority chip}}}}}}}
%    \end{macrocode}
%
% \paragraph{Break-layout headers and shared task cards.}
% The page-breakable renderer needs a header box whose height is synchronized
% across its independent columns. All renderers need one full-width card whose
% surface and text follow the active mode and whose colored west rule identifies
% the target column. The shared card composer places task text first, then an
% optional priority/completion/author row, and finally an optional deadline row.
%
% |\newtcolorbox| defines an environment backed by |tcolorbox|. The header
% accepts two arguments for background and foreground colors. |enhanced| enables
% advanced drawing, |sharp corners| removes rounding, |fontupper| applies bold
% ragged-right text, and |equal height group| gives every header in the current
% board group the height of the tallest member. |valign=center| vertically
% centers shorter headings in that shared height. The card accepts its column
% color as one argument. |width=\linewidth| makes it fill the column;
% |borderline west| draws the semantic accent; the mode macros supply its
% surface, frame, and text; and |kanban/card| is applied last.
%
% |\kanban@rendercard| is a seven-argument deferred composer. |\notblank| from
% |etoolbox| executes metadata rows only when their concatenated fields contain
% material. |\par| begins a new row, |\vspace| separates it, and |\noindent|
% removes paragraph indentation. |\hfill| pushes the author chip to the right;
% a conditional |\hspace| separates adjacent priority and completion pills.
% The same composer is stored for both board layouts, which guarantees
% identical card semantics.
%    \begin{macrocode}
\newtcolorbox{kanban@header}[2]{enhanced,sharp corners,colback=#1,
  boxrule=0pt,coltext=#2,
  fontupper=\bfseries\raggedright,
  left=2mm,right=2mm,
  top=1.2mm,bottom=1.2mm,before skip=0pt,after skip=1.8mm,
  equal height group={KBH@\kanban@boardgroup},valign=center,
  kanban/header}
\newtcolorbox{kanban@card}[1]{enhanced,arc=.6mm,outer arc=.6mm,
  colback=\kanban@cardcolor,colframe=\kanban@bordercolor,
  coltext=\kanban@textcolor,boxrule=.4pt,width=\linewidth,
  borderline west={2.2pt}{0pt}{#1},left=2.4mm,right=2mm,top=1.6mm,
  bottom=1.6mm,before skip=0pt,after skip=1.8mm,
  fontupper=\scriptsize\raggedright,kanban/card}
\newcommand\kanban@rendercard[7]{%
  \begin{kanban@card}{#1}#2%
    \notblank{#3#4#7}{\par\vspace{1mm}\noindent
      \kanban@prioritypill{#3}%
      \notblank{#3}{\notblank{#7}{\hspace{1mm}}{}}{}%
      \kanban@completedchip{#7}\hfill\kanban@authorchip{#4}}{}%
    \notblank{#5}{\par\vspace{1mm}\noindent
      \kanban@deadlinechip{#6}{#5}}{}%
  \end{kanban@card}}
%    \end{macrocode}
%
% \paragraph{Fixed-board containers and renderer adapters.}
% The fixed renderer uses a two-row raster: the first row contains synchronized
% headings and the second contains equal-height transparent column bodies. The
% breakable renderer instead places each header and task list in a separate
% |paracol| stream. Three small adapter commands insert the dynamically stored
% column name, count, color, foreground, and deferred card list into the
% appropriate container.
%
% |kanban@fixedheader| is a |tcolorbox| environment without outer vertical
% skip; raster row equalization supplies its height. |kanban@fixedbody| hides
% the frame, uses |opacityback=0| for a transparent background, removes lateral
% padding, and draws only |borderline south| in a tint of the column color.
% In the adapters, |\csuse| from |etoolbox| constructs and expands a control
% sequence from character data, such as the color for numeric column |#1| or
% the count belonging to the active board group. The break adapter uses the
% equal-height |kanban@header| and emits the stored card list directly after it;
% its final |\par| closes the column material cleanly.
%    \begin{macrocode}
\newtcolorbox{kanban@fixedheader}[2]{enhanced,sharp corners,
  colback=#1,boxrule=0pt,coltext=#2,
  fontupper=\bfseries\raggedright,boxsep=0pt,left=2mm,right=2mm,
  top=1mm,bottom=1mm,before skip=0pt,after skip=0pt,valign=center,
  kanban/header}
\newtcolorbox{kanban@fixedbody}[1]{enhanced,sharp corners,frame hidden,
  colback=white,opacityback=0,boxrule=0pt,
  borderline south={.6pt}{0pt}{#1!55},boxsep=0pt,
  left=0pt,right=0pt,top=2mm,bottom=2mm,
  before skip=0pt,after skip=0pt,valign=top,kanban/column-body}
\newcommand\kanban@onehead[1]{%
  \begin{kanban@fixedheader}{\csuse{kanban@color@#1}}%
      {\csuse{kanban@headingtext@#1}}%
    \csuse{kanban@name@#1}\hfill
      \csuse{kanban@count@\kanban@boardgroup @#1}%
  \end{kanban@fixedheader}}
\newcommand\kanban@onecol[1]{%
  \begin{kanban@fixedbody}{\csuse{kanban@color@#1}}%
    \csuse{kanban@list@\kanban@boardgroup @#1}%
  \end{kanban@fixedbody}}
\newcommand\kanban@onecol@b[1]{%
  \begin{kanban@header}{\csuse{kanban@color@#1}}%
      {\csuse{kanban@headingtext@#1}}%
    \csuse{kanban@name@#1}\hfill
      \csuse{kanban@count@\kanban@boardgroup @#1}%
  \end{kanban@header}%
  \csuse{kanban@list@\kanban@boardgroup @#1}\par}
%    \end{macrocode}
%
%
%    \begin{macrocode}
%</package>
%    \end{macrocode}
